
\section{AFC Stokes seed}\label{sec:afc-basis}

\subsection{Null potential}\label{subsec:afc-null-potential}
\aodprovenance{AFC provenance used here: Null potential; declared distinctions; induced relations; regions; boundaries; Stokes identity on declared cuts.}

\subsection{Stokes identity}\label{subsec:afc-stokes-identity}
The title names the hidden temporal dynamics of Stokes. Use the AFC Stokes identity on declared regions and boundaries:
\begin{equation}
\sum_{i\in R}\operatorname{div}(i;B)
=
\sum_{\substack{i\in R\\j\notin R}}\phi_{ij}(B).
\label{eq:afc-stokes-used}
\end{equation}
A\(\Omega\)D resolves the scoped flux value by
\begin{equation}
\phi_{ij,\AO}(B)=p^D_{ij,\AO}(B)A_{ij,\AO}(B).
\label{eq:afc-flux-ao-resolution}
\end{equation}

\aodremark{The Art of the Leprechaun is the tracing and exploration of Stokes temporal wave dynamics through retained branch/hinge/branch forms. It names the A\(\Omega\)D boundary-accounting stance: follow the cut, name the recurrent forms, and use the finite calculus for scaled prediction entries.}

\subsection{Boundary accounting}\label{subsec:afc-boundary-accounting}
Boundary accounting records current and closure across declared structures. It may be written
\[
\mathbf a\vdash(\mathcal F,\partial\mathcal F,\omega)
\]
when the retained address matters. Inside that context, use ordinary local symbols. The accounting operation records in-flow, exoshedding, reflected return, pending return, shell reclosure, pressure, and SADAR coupling.


\subsection{Radical-free arithmetic policy}\label{subsec:radical-free-arithmetic-policy}
A\(\Omega\)D active arithmetic uses declared finite sums, products, signs, ratios, maxima/minima, and finite powers. Root or radical reconstruction belongs to quarantined external comparison, not to active A\(\Omega\) notation. External comparison maps remain quarantined in the manual.

\subsection{Fractal address}\label{subsec:afc-fractal-address}
The active fractal address is octal. A retained address is written
\[
\mathbf a\in\mathcal A^{\mathrm{ret}}_8,
\qquad
\mathcal A^{\mathrm{ret}}_8\subseteq \{0,1,2,3,4,5,6,7\}^{<\omega}.
\]
The origin anchor is \(\mathtt{00}_8\). The displayed octal ladder begins
\[
\mathtt{00}_8,\ \mathtt{01}_8,\ \mathtt{02}_8,\ \mathtt{03}_8,\ \mathtt{04}_8,\ \mathtt{05}_8,\ \mathtt{06}_8,\ \mathtt{07}_8,\ \mathtt{10}_8,\ldots .
\]
The octal address selects the retained coordinate. The curling-curl support form names the retained cycle form at that coordinate. The A\(\Omega\)D glyph \(a\mid-\) is the fractal coordinate. In particular, \(0\mid-\) binds the statement at anchor \(0\).

\aodprovenance{\eqnref{eq:null-axiom}; \eqnref{eq:null-presentation}; \eqnref{eq:cut-running-form}; \appref{app:fractal-address-bip-biz}.}
